%%%%%%%%%%%%%%%%%%%%%%%%%%%%%%%%%%%%%%%%%%%%%%%%%%%%%%%%%%%%%%%%%%%%%%%%%%%%%%%
%   PEŁNA INSTRUKCJA ROZWIĄZANIA ZADANIA 4 - KANONICZNY PARSER LR(1)
%%%%%%%%%%%%%%%%%%%%%%%%%%%%%%%%%%%%%%%%%%%%%%%%%%%%%%%%%%%%%%%%%%%%%%%%%%%%%%%
\documentclass[11pt, a4paper]{article}

\usepackage[utf8]{inputenc}
\usepackage[T1]{fontenc}
\usepackage{lmodern}
\usepackage[polish]{babel}
\usepackage{geometry}
\geometry{a4paper, margin=2.5cm}

\usepackage{amsmath}
\usepackage{amssymb}
\usepackage{booktabs}
\usepackage{multirow}
\usepackage{float}

\usepackage[unicode=true, colorlinks=true, linkcolor=black, pdfusetitle]{hyperref}

% Definicje komend
\newcommand{\code}[1]{\texttt{#1}}
\newcommand{\taskseparator}{\vspace{0.5cm}\leavevmode\leaders\hrule height 0.5pt\hfill\kern0pt\vspace{0.5cm}}

\title{\textbf{Sprawozdanie z Języków Formalnych i Technik Translacji} \\ Zadanie 4: Konstrukcja pełnego analizatora LR(1)}
\author{Marcel Musiałek \\ Indeks: 279704}
\date{28 stycznia 2026}

\begin{document}

\maketitle
\tableofcontents
\newpage

\section{Treść Zadania}
Celem zadania jest policzenie pełnej tablicy kanonicznego analizatora LR(1) dla podanej gramatyki wyrażeń logicznych:
\begin{align*}
    E &\to E \text{ or } T \mid T \\
    T &\to T \text{ and } F \mid F \\
    F &\to \text{not } F \mid (E) \mid \text{true} \mid \text{false}
\end{align*}

\section{Opis Metody: Analiza LR(1)}

\subsection{Istota kanonicznego lookahead}
Analizator LR(1) różni się od SLR tym, że każda sytuacja wewnątrz stanu (item) posiada przypisany własny zbiór symboli podglądu (\textit{lookahead set}). Sytuacja ma postać $[A \to \alpha \cdot \beta, \{a, b\}]$, gdzie $\{a, b\}$ to terminale, po których wystąpieniu dopuszczalna jest redukcja do symbolu $A$.

\subsection{Algorytm domknięcia (Closure)}
Dla każdej sytuacji $[A \to \alpha \cdot B \beta, L]$ w stanie, dla każdej produkcji $B \to \gamma$, dodajemy sytuację:
\begin{equation}
    [B \to \cdot \gamma, FIRST(\beta L)]
\end{equation}
Dzięki temu lookahead jest precyzyjnie „dziedziczony” z kontekstu, co pozwala uniknąć konfliktów w gramatykach, w których zbiory $FOLLOW$ są zbyt szerokie.

\taskseparator

\section{Pełne Rozwiązanie: Kolekcja Kanoniczna Stanów}

Poniżej przedstawiono pełne stany automatu, z jawnym wyszczególnieniem zbiorów lookahead.

\begin{itemize}
    \item \textbf{I0 (Stan początkowy):} \\
    $[E' \to \cdot E, \{ \$ \}]$ \\
    $[E \to \cdot E \text{ or } T, \{ \$, \text{or} \}]$ \\
    $[E \to \cdot T, \{ \$, \text{or} \}]$ \\
    $[T \to \cdot T \text{ and } F, \{ \$, \text{or}, \text{and} \}]$ \\
    $[T \to \cdot F, \{ \$, \text{or}, \text{and} \}]$ \\
    $[F \to \cdot \text{not } F, \{ \$, \text{or}, \text{and} \}]$ \\
    $[F \to \cdot (E), \{ \$, \text{or}, \text{and} \}]$ \\
    $[F \to \cdot \text{true}, \{ \$, \text{or}, \text{and} \}]$ \\
    $[F \to \cdot \text{false}, \{ \$, \text{or}, \text{and} \}]$

    \item \textbf{I1 ($GOTO(I0, E)$):} \\
    $[E' \to E \cdot, \{ \$ \}]$ $\rightarrow$ \textbf{ACCEPT} \\
    $[E \to E \cdot \text{or } T, \{ \$, \text{or} \}]$

    \item \textbf{I2 ($GOTO(I0, T)$):} \\
    $[E \to T \cdot, \{ \$, \text{or} \}]$ $\rightarrow$ \textbf{Redukcja (2)} \\
    $[T \to T \cdot \text{and } F, \{ \$, \text{or}, \text{and} \}]$

    \item \textbf{I4 ($GOTO(I0, '(')$):} \\
    $[F \to ( \cdot E ), \{ \$, \text{or}, \text{and} \}]$ \\
    $[E \to \cdot E \text{ or } T, \{ ), \text{or} \}]$ \\
    $[E \to \cdot T, \{ ), \text{or} \}]$ \\
    $[T \to \cdot T \text{ and } F, \{ ), \text{or}, \text{and} \}]$ \\
    $[T \to \cdot F, \{ ), \text{or}, \text{and} \}]$ \\
    (Zauważmy zmianę symbolu $\$$ na $)$ w zbiorze podglądu).

    \item \textbf{I10 ($GOTO(I3, F)$):} \\
    $[F \to \text{not } F \cdot, \{ \$, \text{or}, \text{and} \}]$ $\rightarrow$ \textbf{Redukcja (5)}
\end{itemize}

\taskseparator

\section{Pełna Tabela Analizatora LR(1)}

Tabela zawiera 12 kolumn, uwzględniając wszystkie terminale (ACTION) i nieterminale (GOTO).


\begin{table}[H]
\centering
\scriptsize
\begin{tabular}{|c|cccccccc|ccc|}
\hline
\multirow{2}{*}{\textbf{Stan}} & \multicolumn{8}{c|}{\textbf{ACTION}} & \multicolumn{3}{c|}{\textbf{GOTO}} \\ \cline{2-12} 
 & \textbf{or} & \textbf{and} & \textbf{not} & \textbf{(} & \textbf{)} & \textbf{true} & \textbf{false} & \textbf{\$} & \textbf{E} & \textbf{T} & \textbf{F} \\ \hline
0 & & & s3 & s4 & & s8 & s9 & & 1 & 2 & 5 \\
1 & s6 & & & & & & & acc & & & \\
2 & r2 & s7 & & & & & & r2 & & & \\
3 & & & s3 & s4 & & s8 & s9 & & & & 10 \\
4 & & & s14 & s15 & & s25 & s26 & & 11 & 12 & 13 \\
5 & r4 & r4 & & & & & & r4 & & & \\
6 & & & s3 & s4 & & s8 & s9 & & & 16 & 5 \\
7 & & & s3 & s4 & & s8 & s9 & & & & 17 \\
8 & r7 & r7 & & & & & & r7 & & & \\
9 & r8 & r8 & & & & & & r8 & & & \\
10 & r5 & r5 & & & & & & r5 & & & \\
11 & s19 & & & & s18 & & & & & & \\
12 & r2 & s20 & & & r2 & & & & & & \\
13 & r4 & r4 & & & r4 & & & & & & \\
14 & & & s14 & s15 & & s25 & s26 & & & & 21 \\
15 & & & s14 & s15 & & s25 & s26 & & 11 & 12 & 13 \\
16 & r1 & s7 & & & & & & r1 & & & \\
17 & r3 & r3 & & & & & & r3 & & & \\
18 & r6 & r6 & & & & & & r6 & & & \\
19 & & & s14 & s15 & & s25 & s26 & & & 23 & 13 \\
20 & & & s14 & s15 & & s25 & s26 & & & & 24 \\
21 & r5 & r5 & & & r5 & & & & & & \\
22 & r6 & r6 & & & r6 & & & & & & \\
23 & r1 & s20 & & & r1 & & & & & & \\
24 & r3 & r3 & & & r3 & & & & & & \\
25 & r7 & r7 & & & r7 & & & & & & \\
26 & r8 & r8 & & & r8 & & & & & & \\ \hline
\end{tabular}
\end{table}

\section{Wnioski}
Kanoniczny analizator LR(1) wyeliminował niejednoznaczności poprzez precyzyjne rozdzielenie stanów o tych samych rdzeniach, ale różnych kontekstach (np. stany I10 i I21). Jest to metoda najbardziej ogólna, choć kosztowna pod względem liczby stanów w porównaniu do metod SLR.

\end{document}